\chapter{Conclusion and Future Work}
\label{ch:conclusion}

\section{Conclusion}

This project set out to build a reliable, reproducible and fully automated
pipeline for segmenting glioblastoma sub-regions from multi-parametric brain MRI.
Using the BraTS~2020 dataset restricted to its 293 high-grade-glioma cases, we
trained a 3D~U-Net that takes the four co-registered MRI modalities as input and
produces a dense voxel-wise segmentation into necrotic/non-enhancing core,
peritumoral edema and enhancing tumour. The model was optimised with a combined
Dice and cross-entropy loss and evaluated with both an overlap metric (Dice) and a
boundary metric (HD95). It reaches a foreground mean Dice of 0.71 on the held-out
validation split, localising the tumour reliably and segmenting the edema and
enhancing tumour well, with the expected drop in accuracy on the smaller and more
ambiguous necrotic core. On the standard composite BraTS regions the model's
whole-tumour Dice (0.88) sits within a few points of the top BraTS~2020
challenge entries, despite training for only 34 epochs on a single 8~GB GPU.
Just as important, the whole study is reproducible: fixed seeds, a
deterministic split, a single configurable entry point and a notebook that
runs the pipeline end to end on a single consumer GPU. The trained model is
also deployed behind a public, interactive web application
(\url{https://glioblastome-brain-tumor-detection.streamlit.app/}), where a user
can select an example scan or upload their own and receive a tumour detection
and segmentation result directly, turning the research pipeline into a usable
tool rather than a static report.

\section{Future work}

Several concrete extensions follow directly from the limitations in
Chapter~\ref{ch:discussion}:

\begin{itemize}
  \item \textbf{Stronger encoders.} The pipeline is architecture-agnostic;
        swapping the 3D~U-Net for an attention-gated variant
        \citep{oktay2018attentionunet} or a transformer encoder such as
        Swin~UNETR \citep{hatamizadeh2022swinunetr} is a natural next step now
        that the data and training loop are in place.
  \item \textbf{Whole-volume inference.} Replacing the fixed centre-crop with
        sliding-window inference would remove any risk of clipping off-centre
        tumours and would align evaluation with standard BraTS practice.
  \item \textbf{Broader data and external validation.} Training on the full BraTS
        set (including LGG) and testing on an independent cohort would establish
        how well the model generalises across grades, scanners and institutions.
  \item \textbf{Uncertainty estimation.} Adding test-time dropout or ensembling
        would let the model express confidence in its predictions, which is
        valuable in a clinical decision-support setting.
  \item \textbf{Extending the deployed application.} The web application
        described above currently requires all four modalities in BraTS-style
        pre-processed form; supporting raw, non-registered clinical exports
        would need automatic skull-stripping and co-registration added to the
        upload path, making the tool usable directly on unprocessed scans
        rather than only on research-format data.
\end{itemize}

Taken together, this work delivers a correct and honestly evaluated glioblastoma
segmentation pipeline and a clear path from that pipeline towards a deployable
tool.
